% GENERATED FILE. Edit canonical lessons, then run npm run generate:books.
% canonical-source-hash: fnv1a64:7adb0f14449f1dcd
% canonical-lessons: LA-C374-noceo, LA-C374-prosum, LA-C374-obsum, LA-C374-impedio, LA-C374-perturbo

\chapter{Harm and Help}
\label{ch:la-harm-and-help}
\hlchaptermodality{374}

\begin{chapteropening}
\textbf{By the end of this chapter:} \emph{I can say that something harms or helps, does damage, stands in the way of a plan or upsets someone.}

Complete the last lesson of chapter 374: I can say that something harms or helps, does damage, stands in the way of a plan or upsets someone.
\end{chapteropening}

\section[noceō, nocēre]{noceō, nocēre — to harm, to hurt}

\label{lesson:LA-C374-noceo}



\begin{quote}
Before the new one: say the Latin for to turn out well, then the Latin for to overflow.
\end{quote}

\subsection*{You'll want to know: noceō, nocēre}
\textbf{noceō, nocēre} — \textquotedblleft{}to harm, to hurt\textquotedblright{}.

\begin{quote}
Fūmus oculīs nocet.
\end{quote}

\begin{quote}
\emph{Smoke harms the eyes.}
\end{quote}

\emph{noceō, nocēre, nocuī, nocitum} + dative: \emph{tibi nocet}, it harms you. English \emph{innocent} (not harming) and \emph{nuisance} come from its family.

\begin{cousinweb}
\emph{noceō} is inherited from Proto-Italic \emph{*nokeō} (EtymDB), a causative of the root of \emph{nex} \textquotedblleft{}violent death\textquotedblright{}: to make perish, then to harm. English \emph{innocent} and \emph{nuisance} come from its family.
\end{cousinweb}

\begin{grammarlens}[title={What you've built}]
One more word for telling, explaining and describing at B1.
\end{grammarlens}

\subsection*{Your turn}
\begin{itemize}
\raggedright
  \item \emph{Say it:} \emph{noceō, nocēre}
  \item \emph{Say it:} \emph{noceō, nocēre}, once more
  \item \emph{Say it:} a sentence of your own with \emph{noceō, nocēre}
  \item \emph{Recall:} say the Latin for to turn out well, then the Latin for to overflow, then say \emph{noceō, nocēre} again
\end{itemize}

\begin{tcolorbox}[breakable,colback=teal!4,colframe=teal!35!black,title={Before you move on}]
What does noceō, nocēre mean? (\textquotedblleft{}To harm, to hurt\textquotedblright{}.) Say it once more.
\end{tcolorbox}
\section[prōsum, prōdesse]{prōsum, prōdesse — to be useful, to benefit}

\label{lesson:LA-C374-prosum}



\begin{quote}
Before the new one: say the Latin for to overflow, then the Latin for to harm.
\end{quote}

\subsection*{You'll want to know: prōsum, prōdesse}
\textbf{prōsum, prōdesse} — \textquotedblleft{}to be useful, to benefit\textquotedblright{}.

\begin{quote}
Haec medicīna tibi prōderit.
\end{quote}

\begin{quote}
\emph{This medicine will do you good.}
\end{quote}

\emph{prōsum, prōdesse, prōfuī} + dative. It goes like \emph{sum}, with \emph{prōd-} before a vowel: \emph{prōdest}, \emph{prōderit}.

\begin{cousinweb}
\emph{prōsum} is built from \emph{prō} \textquotedblleft{}for, on behalf of\textquotedblright{} and \emph{sum} \textquotedblleft{}to be\textquotedblright{} (EtymDB; Lewis \& Short), with an old \emph{-d-} kept before a vowel: \emph{prōdest}.
\end{cousinweb}

\begin{grammarlens}[title={What you've built}]
One more word for telling, explaining and describing at B1.
\end{grammarlens}

\subsection*{Your turn}
\begin{itemize}
\raggedright
  \item \emph{Say it:} \emph{prōsum, prōdesse}
  \item \emph{Say it:} \emph{prōsum, prōdesse}, once more
  \item \emph{Say it:} a sentence of your own with \emph{prōsum, prōdesse}
  \item \emph{Recall:} say the Latin for to overflow, then the Latin for to harm, then say \emph{prōsum, prōdesse} again
\end{itemize}

\begin{tcolorbox}[breakable,colback=teal!4,colframe=teal!35!black,title={Before you move on}]
What does prōsum, prōdesse mean? (\textquotedblleft{}To be useful, to benefit\textquotedblright{}.) Say it once more.
\end{tcolorbox}
\section[obsum, obesse]{obsum, obesse — to be a hindrance, to do harm}

\label{lesson:LA-C374-obsum}



\begin{quote}
Before the new one: say the Latin for to harm, then the Latin for to be useful.
\end{quote}

\subsection*{You'll want to know: obsum, obesse}
\textbf{obsum, obesse} — \textquotedblleft{}to be a hindrance, to do harm\textquotedblright{}.

\begin{quote}
Nimia pecūnia saepe obest.
\end{quote}

\begin{quote}
\emph{Too much money often does harm.}
\end{quote}

\emph{obsum, obesse, obfuī} + dative: the opposite of \emph{prōsum}. \emph{Nōn obest}, it does no harm.

\begin{cousinweb}
\emph{obsum} is built from \emph{ob} \textquotedblleft{}against, in the way\textquotedblright{} and \emph{sum} \textquotedblleft{}to be\textquotedblright{} (EtymDB; Lewis \& Short): to be in someone's way.
\end{cousinweb}

\begin{grammarlens}[title={What you've built}]
One more word for telling, explaining and describing at B1.
\end{grammarlens}

\subsection*{Your turn}
\begin{itemize}
\raggedright
  \item \emph{Say it:} \emph{obsum, obesse}
  \item \emph{Say it:} \emph{obsum, obesse}, once more
  \item \emph{Say it:} a sentence of your own with \emph{obsum, obesse}
  \item \emph{Recall:} say the Latin for to harm, then the Latin for to be useful, then say \emph{obsum, obesse} again
\end{itemize}

\begin{tcolorbox}[breakable,colback=teal!4,colframe=teal!35!black,title={Before you move on}]
What does obsum, obesse mean? (\textquotedblleft{}To be a hindrance, to do harm\textquotedblright{}.) Say it once more.
\end{tcolorbox}
\section[impediō, impedīre]{impediō, impedīre — to hinder, to hamper}

\label{lesson:LA-C374-impedio}



\begin{quote}
Before the new one: say the Latin for to be useful, then the Latin for to be a hindrance.
\end{quote}

\subsection*{You'll want to know: impediō, impedīre}
\textbf{impediō, impedīre} — \textquotedblleft{}to hinder, to hamper\textquotedblright{}.

\begin{quote}
Nix iter nostrum impedīvit.
\end{quote}

\begin{quote}
\emph{Snow hampered our journey.}
\end{quote}

\emph{impediō, impedīre, impedīvī, impedītum}. Built on \emph{pēs, pedis}: to shackle the feet. \emph{Impediō quōminus} + subjunctive, to keep from. English \emph{impede} and \emph{impediment}.

\begin{cousinweb}
\emph{impediō} is built from \emph{in} and \emph{pēs, pedis} \textquotedblleft{}a foot\textquotedblright{} (Lewis \& Short): to catch someone by the feet. English \emph{impede} and \emph{impediment} come from it.
\end{cousinweb}

\begin{grammarlens}[title={What you've built}]
One more word for telling, explaining and describing at B1.
\end{grammarlens}

\subsection*{Your turn}
\begin{itemize}
\raggedright
  \item \emph{Say it:} \emph{impediō, impedīre}
  \item \emph{Say it:} \emph{impediō, impedīre}, once more
  \item \emph{Say it:} a sentence of your own with \emph{impediō, impedīre}
  \item \emph{Recall:} say the Latin for to be useful, then the Latin for to be a hindrance, then say \emph{impediō, impedīre} again
\end{itemize}

\begin{tcolorbox}[breakable,colback=teal!4,colframe=teal!35!black,title={Before you move on}]
What does impediō, impedīre mean? (\textquotedblleft{}To hinder, to hamper\textquotedblright{}.) Say it once more.
\end{tcolorbox}
\section[perturbō, perturbāre]{perturbō, perturbāre — to throw into confusion, to upset}

\label{lesson:LA-C374-perturbo}



\begin{quote}
Before the new one: say the Latin for to be a hindrance, then the Latin for to hinder.
\end{quote}

\subsection*{You'll want to know: perturbō, perturbāre}
\textbf{perturbō, perturbāre} — \textquotedblleft{}to throw into confusion, to upset\textquotedblright{}.

\begin{quote}
Nūntius malus mē perturbāvit.
\end{quote}

\begin{quote}
\emph{The bad news upset me.}
\end{quote}

\emph{perturbō, perturbāre, perturbāvī, perturbātum}. \emph{Per-}, thoroughly, with a verb for stirring up. English \emph{perturb}.

\begin{cousinweb}
\emph{perturbō} is built from \emph{per} \textquotedblleft{}thoroughly\textquotedblright{} and \emph{turbō} \textquotedblleft{}to stir up, to disturb\textquotedblright{} (EtymDB; Lewis \& Short), which belongs with \emph{turba} \textquotedblleft{}an uproar\textquotedblright{}. English \emph{perturb} comes from it.
\end{cousinweb}

\begin{grammarlens}[title={What you've built}]
One more word for telling, explaining and describing at B1.
\end{grammarlens}

\subsection*{Your turn}
\begin{itemize}
\raggedright
  \item \emph{Say it:} \emph{perturbō, perturbāre}
  \item \emph{Say it:} \emph{perturbō, perturbāre}, once more
  \item \emph{Say it:} a sentence of your own with \emph{perturbō, perturbāre}
  \item \emph{Recall:} say the Latin for to be a hindrance, then the Latin for to hinder, then say \emph{perturbō, perturbāre} again
\end{itemize}

\begin{tcolorbox}[breakable,colback=teal!4,colframe=teal!35!black,title={Before you move on}]
What does perturbō, perturbāre mean? (\textquotedblleft{}To throw into confusion, to upset\textquotedblright{}.) Say it once more.
\end{tcolorbox}
